%%% ============================================================
%%% PAPER 1 SKELETON: The Freeze-Thaw Transit
%%% A Dual-Regime Scalar Dark-Energy Model with an
%%% Analytic Turnaround at e^{-1}
%%%
%%% B.R. Sanders, M. Gish (2026)
%%% Target venue: JCAP
%%%
%%% STATUS: SKELETON for ClaudeCode expansion.
%%% Source materials for fill-in:
%%%   - Existing detailed draft (logarithmic_quintessence_v3.tex)
%%%   - desi_xi_optimize_v2.py (the actual fit script)
%%%   - paper1_patches.tex (§6.X Type-T/F/A subsection, F5 criterion)
%%% Companion: paper2_substrate_M22_skeleton.tex
%%%
%%% UNIT NOTE FOR EXPANSION: the supplementary script v2 uses
%%%   N_start = -4  ==> z_init ~ 54 (NOT z=20 as text says)
%%%   xi_dot     ==> dxi/dN (e-folding derivative)
%%%   Lambda4    ==> Λ^4/(3 H_0^2 M_Pl^2)  [= 0.231 in fit]
%%% Manuscript text "z_i ≈ 20" is a discrepancy to harmonize.
%%% ============================================================

\documentclass[11pt,reqno]{amsart}

\usepackage{amsmath,amssymb,amsthm,mathtools}
\usepackage[margin=1in]{geometry}
\usepackage[colorlinks=true,linkcolor=blue,citecolor=blue,urlcolor=blue]{hyperref}
\usepackage{microtype}
\usepackage{array}

%% theorem environments
\theoremstyle{plain}
\newtheorem{theorem}{Theorem}
\newtheorem{proposition}[theorem]{Proposition}
\newtheorem{corollary}[theorem]{Corollary}
\newtheorem{lemma}[theorem]{Lemma}

\theoremstyle{definition}
\newtheorem{definition}[theorem]{Definition}

\theoremstyle{remark}
\newtheorem*{remark}{Remark}

%% macros
\newcommand{\Pl}{\mathrm{Pl}}
\newcommand{\SM}{\mathrm{SM}}
\newcommand{\Gibbs}{\mathrm{Gibbs}}
\newcommand{\DE}{\mathrm{DE}}
\DeclareMathOperator{\diag}{diag}

%% TODO macro for ClaudeCode expansion targets
\newcommand{\TODO}[1]{\par\noindent\fbox{\parbox{0.95\linewidth}{\small\textbf{TODO:} #1}}\par}

%% title matter
\title[The Freeze-Thaw Transit]{The Freeze-Thaw Transit:\\
A Dual-Regime Scalar Dark-Energy Model with\\
an Analytic Turnaround at $e^{-1}$}

\author{B.~R.~Sanders}
\address{7Site LLC, Hot Springs, Arkansas, USA}
\email{brayden@7site.co}

\author{M.~Gish}
\address{Independent Researcher, Hot Springs, Arkansas, USA}
\email{monica.gish1992@gmail.com}

\date{\today}

\keywords{quintessence, dark energy, freeze-thaw transit, scalar field,
logarithmic potential, analytic vacuum, paradox classification, Stage-IV
cosmology}

\subjclass[2020]{83F05, 83C56, 85A40}

\begin{document}

\begin{abstract}
Scalar-field models of dark energy traditionally fall into two
dynamical classes: \emph{freezing quintessence}, where $w(z)$
approaches $-1$ from above asymptotically, and \emph{thawing
quintessence}, where $w(z)$ departs from $-1$ over the observable
epoch. A given quintessence potential typically produces only one of
these regimes within its physical history. We study a minimally
coupled real, positive, dimensionless scalar field $\Xi$ with
self-interaction $V(\Xi) = \Lambda^{4}\,\Xi\log\Xi$, where $\Lambda$
is a mass scale that sets the dark-energy density. The potential
admits an analytic minimum at $\Xi_{0} = e^{-1}$, independent of the
coupling. Linearizing about the vacuum gives a stable massive scalar
with $m_\Xi^{2} = \Lambda^{4} e / M_\Pl^{2}$, numerically near the
Hubble scale today when $\Lambda \approx 1.7$~meV.

We show that the model's FRW trajectory exhibits \emph{both freezing
and thawing within a single physical history}: an outbound thaw at
high redshift, a momentary frozen turnaround at $z \approx 2$ in the
documented fit (where $\dot\Xi = 0$ and $w_\Xi = -1$ instantaneously),
and an asymptotic refreeze toward $\Xi_{0}$ in the cosmological future.
The current cosmological epoch lies in the post-turnaround thaw phase.
We classify the trajectory using a three-regime taxonomy (T-thaw,
F-frozen, A-asymptotic) and identify the vacuum $\Xi_{0} = e^{-1}$ as
the geometric anchor of the freeze-thaw transit, simultaneously the
minimum of $V$ and the maximizer of the per-bin Gibbs functional
$H_\Gibbs(x) = -x\log x$ on $(0, 1]$, a structural correspondence we
discuss as motivational rather than derivational.

A DESI~2024 DR1 $(w_0, w_a)$ consistency check places the documented
fit within $\sim 1.2\sigma$ of the published Gaussian summary under
the uncorrelated approximation; under the correlated metric
($\rho \approx -0.9$) $\Lambda$CDM remains marginally preferred. The
$\chi^{2}$ values reported quantify proximity to the published
$(w_0, w_a)$ marginal Gaussian and are not derived from the
underlying joint BAO + CMB + SN likelihood, which is deferred to a
companion numerical paper. Five falsifiability criteria $(F_1$--$F_5)$
are testable at Stage-IV precision, including a new criterion specific
to the dual-regime trajectory: a non-monotone $w(z)$ with a local
minimum near $-1$ at intermediate redshift. We position this paper as
the dynamical-setup component of a synthesis program on
substrate-algebraic dark-energy structure; a companion paper
\cite{SandersGish2026M22} derives $\Omega_\DE = 479/700 = 0.6843$
(matching Planck~2018 to $0.06\%$) from substrate algebra mediated by
the Mathieu group $M_{22}$.
\end{abstract}

\maketitle

%%% ============================================================
\section{Introduction}\label{sec:introduction}

\TODO{Open with DESI~2024 DR1 / DR2 results
(\cite{DESI2024BAO,DESI2024VI,DESI2025DR2}) suggesting dynamical
dark energy at $2.8$--$4.2\sigma$ in joint analyses; review revival
of quintessence \cite{RatraPeebles1988,Wetterich1988,Frieman1995}.}

\subsection{The freezing/thawing dichotomy}

\TODO{Present the Caldwell--Linder \cite{CaldwellLinder2005}
classification of quintessence into freezing vs. thawing regimes;
note that standard potentials produce only one regime per trajectory;
cite examples (inverse power law $\to$ freezing; PNGB $\to$ thawing).
Argue that the observational picture (averaged $w \approx -1$ in the
recent past, but DESI-preferred dynamical $(w_0, w_a)$ at $z \lesssim
1.5$) is naturally read as both regimes occurring in one history,
not as the universe selecting between them.}

\subsection{The minimal model}

\TODO{Introduce $V(\Xi) = \Lambda^{4} \Xi \log \Xi$ with dimensionless
$\Xi$. State the two structural anchors: analytic vacuum at
$\Xi_{0} = e^{-1}$, fluctuation mass
$m_\Xi^{2} = \Lambda^{4} e / M_\Pl^{2}$. Note structural rhyme with
Gibbs functional $H_\Gibbs(x) = -x \log x$ and reference
\S\ref{sec:entropy} for the explicit caveat that this is a structural
correspondence, not a derivation.}

\subsection{The freeze-thaw transit}

\TODO{Preview the Type-T/F/A regime structure: outbound thaw, frozen
turnaround at $z \approx 2$, asymptotic refreeze. Note that the
current epoch is in the post-turnaround Type-T regime. State that
this dual-regime trajectory is the natural reading of the
observational picture.}

\subsection{Scope and the synthesis program}

\TODO{This paper establishes dynamical content (action, vacuum, mass,
FRW evolution, Type-T/F/A classification). $\Lambda$ is fixed
phenomenologically; $\Omega_\DE$ is taken as Planck input. Companion
paper \cite{SandersGish2026M22} derives $\Omega_\DE = 479/700$ from
substrate algebra. Together: dynamical setup (here) + structural
source (companion).}

\subsection{Organization}

\TODO{One-paragraph overview of \S\S 2--10.}

%%% ============================================================
\section{Canonical Action}\label{sec:action}

\subsection{Field content and conventions}\label{sec:fieldcontent}

\TODO{Field content $\{g_{\mu\nu}, \mathrm{SM}, \Xi(x)\}$;
$(-,+,+,+)$ signature; $\Xi$ real, positive, dimensionless, gauge
singlet. Field domain $\Xi > 0$. Footnote on $\psi = \log\Xi$
reparametrization and on canonical $\phi = M_\Pl \Xi$ alternative;
state we work in $\Xi$ for vacuum-transparency.}

\subsection{Action}

\TODO{Equation \eqref{eq:action} — state minimal action with
$M_\Pl^{2}$ kinetic prefactor and $\Lambda^{4}$ potential
prefactor. Verify dimensional analysis. Note single
model-specific parameter $\Lambda$. End with remark that
information-theoretic origin story is supplementary, deferred to
\S\ref{sec:entropy}.}

\begin{equation}\label{eq:action}
\TODO{[INSERT MINIMAL ACTION FROM EXISTING DRAFT]}
\end{equation}

%%% ============================================================
\section{Derivations}\label{sec:derivations}

\subsection{Euler-Lagrange equation}\label{sec:EL}

\TODO{Derive
$M_\Pl^{2}\,\Box\Xi = \Lambda^{4}(1 + \log\Xi)$
(label \texttt{eq:EL}); note the characteristic mass scale
$\Lambda^{4}/M_\Pl^{2}$. Vacuum condition $\Box\Xi = 0$
gives $1 + \log\Xi_{0} = 0$ coupling-independently.}

\subsection{Stress-energy tensor}\label{sec:stress}

\TODO{Derive
$T^\Xi_{\mu\nu}$ from variation w.r.t.\ $g^{\mu\nu}$
(label \texttt{eq:stress}). Confirm covariant conservation on
shell.}

\subsection{Vacuum/ground state}\label{sec:vacrenorm}

\TODO{Derive $\Xi_{0} = e^{-1}$, $V'(\Xi_{0}) = 0$,
$V''(\Xi_{0}) = \Lambda^{4} e > 0$. Discuss bare vs renormalized
vacuum: $\Lambda_\mathrm{obs} = \Lambda_\mathrm{bare} - \Lambda^{4}/e$.
This subsection is critical for the §6.2 energy-budget framing.}

\subsection{Fluctuation mass and stability}\label{sec:mass}

\TODO{Linearize about $\Xi_0$: $m_\Xi^{2} = \Lambda^{4} e / M_\Pl^{2}$
(label \texttt{eq:mass}). Note dark-energy scale match: $\Lambda \sim
1.5$~meV from dimensional analysis, $\Lambda \approx 1.7$~meV from
the documented fit; discrepancy explained by geometric factor
$(3 e \Omega_\Lambda)^{1/4}$.}

\subsection{Global behavior of $V$}\label{sec:globalV}

\TODO{$V \to 0^-$ as $\Xi \to 0^+$; $V(\Xi_0) = -\Lambda^4/e$ global
minimum on $\Xi > 0$; $V \to +\infty$ as $\Xi \to \infty$.}

%%% ============================================================
\section{FRW Cosmology}\label{sec:frw}

\subsection{Background equations}\label{sec:bgequations}

\TODO{In flat FRW: derive $\rho_\Xi$, $p_\Xi$
(labels \texttt{eq:rho}, \texttt{eq:p}). Friedmann equations
\texttt{eq:friedmann}. Field EoM in FRW \texttt{eq:xieom}:
$M_\Pl^{2}(\ddot\Xi + 3H\dot\Xi) = -\Lambda^{4}(1 + \log\Xi)$.}

\subsection{Equation of state}\label{sec:eos}

\TODO{Derive
$w_\Xi = -1 + M_\Pl^{2}\dot\Xi^{2}/\rho_\Xi$ (label \texttt{eq:w}).
State Proposition on limiting regimes:
(i) $w_\Xi = -1$ at vacuum;
(ii) $w_\Xi \to -1$ from above on rolling branch;
(iii) $w_\Xi \to +1$ in kinetic limit.}

\subsection{Late-time attractor}\label{sec:lateattractor}

\TODO{Damping argument: $3H\dot\Xi \to 0$, $\Xi \to \Xi_0$,
$w_\Xi \to -1$. This is the Type-A asymptotic regime per
\S\ref{sec:typeTFA} below.}

\subsection{Type-T/F/A classification of the trajectory}\label{sec:typeTFA}

\TODO{INSERT THE FULL §6.X SUBSECTION FROM \texttt{paper1\_patches.tex}.
Three regimes (T, F, A). Geometric anchor at $\Xi_0 = e^{-1}$.
Distinction from Caldwell--Linder dichotomy. Forward-reference $F_5$.}

%%% ============================================================
\section{Structural Connections to Information-Theoretic Functionals}\label{sec:entropy}

\subsection{The Gibbs functional}

\TODO{$H_\Gibbs(x) = -x \log x$ on $(0, 1]$, max at $x = e^{-1}$.
Formal relation $V(\Xi) = -\Lambda^{4} H_\Gibbs(\Xi)$. Explicit
caveat: $\Xi$ is unbounded scalar, not probability density;
relation is structural, not operational.}

\subsection{Bialynicki-Birula--Mycielski nonlinearity}

\TODO{Cite \cite{BialynickiBirulaMycielski1976}; note BBM uniqueness
applies to $|\psi|^2 \log |\psi|^2$ on probability densities, not to
scalar fields. Structural rhyme, not derivation.}

\subsection{Information-theoretic precedents}

\TODO{Brief note on \cite{Verlinde2011,CHLZ2012,Ensslin2013,Caticha2012}
as adjacent literature; structural rather than derivational connection.
The empirical content of the paper (\S\ref{sec:obs}) does not depend
on these connections.}

%%% ============================================================
\section{Observational Predictions}\label{sec:obs}

\subsection{The $w(z)$ trajectory on the rolling branch}\label{sec:wzrolling}

\TODO{Reframe this subsection from existing draft: the documented fit
shows the post-turnaround thaw on the inbound leg toward $\Xi_0$,
\emph{not} a pure freezing trajectory. Set up the DESI consistency
check that follows.}

\subsection{Numerical integration and DESI consistency check}\label{sec:desifit}

\TODO{The supplementary script \texttt{desi\_xi\_optimize\_v2.py}
performs a $30 \times 25 \times 15 = 11{,}250$-point grid search over
$(\Lambda^4, \Xi_i, \dot\Xi_i)$ at $N_\mathrm{start} = -4$
($z_\mathrm{init} \approx 54$ — note discrepancy with manuscript text
\textbf{``z\_i ≈ 20''}; harmonize on expansion). Best fit at
$\chi^2_{\rho=-0.9}$ minimum. Report:
\begin{itemize}
\item Best-fit parameters $(\Lambda^4, \Xi_i, \dot\Xi_i)$ from script
output
\item $w(z)$ table at $z \in \{0, 0.3, 0.5, 0.8, 1, 1.5, 2\}$
\item $\chi^2_\mathrm{Gauss,\rho=0}$ and $\chi^2_\mathrm{Gauss,\rho=-0.9}$
for $\Xi$ model and $\Lambda$CDM
\item Honest framing: under the correlated metric
($\rho = -0.9$, the disciplined comparison), $\Lambda$CDM remains
marginally preferred by $\Delta\chi^2 \approx 5.6$. Lead with this.
\end{itemize}}

\paragraph{Energy-budget framing.}
\TODO{Rewrite from existing draft. The reported $\Omega_\Xi$
fraction includes both the renormalized cosmological constant and
the scalar's dynamical deviation. Distinguish bare $w_\Xi$ from
combined-sector effective $w_\DE$ tracked by the script.}

\paragraph{On the freeze-thaw transit.}
\TODO{Reframe ``On the freezing label'' from existing draft to
emphasize the dual-regime structure: outbound thaw $\to$ frozen
turnaround at $z \approx 2$ $\to$ inbound thaw $\to$ asymptotic
refreeze. Current epoch is in Type-T post-turnaround regime.}

\paragraph{Initial-condition tuning.}
\TODO{Honest acknowledgment of $(\Xi_i, \dot\Xi_i)$ tuning.
Note that the outbound initial direction $\dot\Xi_i > 0$ is itself
a choice. No naturalness argument singles out these values; the
freezing-branch + dual-regime structure restricts but does not
derive them. Reduction to fewer free parameters via attractor
mechanism is target for follow-up.}

\paragraph{Methodological scope of this consistency check.}
\TODO{Restate explicitly: $\chi^2_\mathrm{Gauss}$ values quantify
proximity to published $(w_0, w_a)$ Gaussian summary, not
goodness-of-fit to underlying joint likelihood. Full DR2 + Pantheon+ +
Planck likelihood is deferred to companion numerical paper.}

\paragraph{Reproducibility.}
\TODO{State that \texttt{desi\_xi\_optimize\_v2.py} reproduces the
fit. Total runtime $\sim 30$~s on standard laptop. Verification
script \texttt{proof\_xi\_canonical.py} runs 22 algebraic and
stability tests. Both archived at DOI \texttt{10.5281/zenodo.18852047}.}

\subsection{Falsification criteria}\label{sec:falsify}

\TODO{Insert the existing F1--F4 from the draft (do not replace).
Then append F5 from \texttt{paper1\_patches.tex}.}

\begin{enumerate}
\item[$(F_1)$] \TODO{$w_\Xi(z) \geq -1$ on rolling branch; phantom
preference at $\geq 3\sigma$ in joint DR2 + Pantheon+ + Planck would
disfavor.}

\item[$(F_2)$] \TODO{Monotone $w(z) \to -1$ from above over
$0 \leq z \leq 1.5$; non-monotone over that range at $\geq 3\sigma$
would falsify the rolling-branch fit.}

\item[$(F_3)$] \TODO{Asymptotic $w \to -1$; late-time inconsistency
at $\geq 3\sigma$ disfavors.}

\item[$(F_4)$] \TODO{Specific $w(z=0.5) = -0.943 \pm 0.05$ at Stage-IV
precision; deviation would disfavor documented fit.}

\item[$(F_5)$] \TODO{INSERT F5 FROM \texttt{paper1\_patches.tex}.
Non-monotone $w(z)$ with local minimum near $-1$ at intermediate
redshift; the dual-regime distinguishing prediction.}
\end{enumerate}

\subsection{Fifth-force and local tests}\label{sec:fifthforce}

\TODO{No direct matter coupling in minimal action; gravitational
strength only; consistent with E\"ot-Wash, MICROSCOPE
\cite{EotWash,MICROSCOPE}.}

\subsection{EFT validity over field excursion}\label{sec:eft}

\TODO{Acknowledge transplanckian field excursion
$\Delta\phi \sim 0.5\,M_\Pl$ in canonical variables. No shift
symmetry. UV completion is open question (shared with most
quintessence). Tree-level effective treatment.}

\subsection{Perturbations}\label{sec:perts}

\TODO{Rest-frame sound speed $c_s^2 = 1$ (canonical scalar).
Standard quintessence subhorizon clustering suppression.
Full perturbation analysis (ISW, CMB lensing) deferred to
companion numerical paper.}

\subsection{Gravitational-wave speed}\label{sec:gw}

\TODO{Tensor perturbations propagate at $c$; consistent with GW170817
\cite{GW170817}. No predicted superluminal signal.}

%%% ============================================================
\section{Comparison to Literature}\label{sec:comparison}

\TODO{Insert comparison table from existing draft. Discuss
\cite{RatraPeebles1988,Wetterich1988,Frieman1995,BarrowParsons1995,
ColemanWeinberg1973,Thompson2019,KouwnOh2012,Chavanis2015,Oikonomou2019}.
Add row for \emph{dual-regime trajectory} characterizing the present
work's distinctive feature. Update novelty remark to current literature
search date.}

%%% ============================================================
\section{Scope and Limitations}\label{sec:scope}

\TODO{Two-sided list: what the paper claims (locked) vs.\ what it does
\emph{not} claim. Match parallel structure to companion paper
\cite{SandersGish2026M22}. Explicit non-claims: derivation of $\Lambda$
from microphysics, derivation of $\Omega_\DE$ (deferred to companion),
direct-matter coupling, mathematical uniqueness theorem.}

%%% ============================================================
\section{The Synthesis Program}\label{sec:synthesis}

\TODO{This paper as the dynamical-setup component of a larger program.
Companion papers in the synthesis program:
\begin{itemize}
\item \cite{SandersGishSigma2026}: Substrate non-associativity decay over $\mathbb{Z}/N\mathbb{Z}$
\item \cite{SandersGishFourCore2026}: Joint closure / four-core attractor on $\mathbb{Z}/10\mathbb{Z}$
\item \cite{SandersGish2026M22}: Substrate-algebraic derivation of $\Omega_\DE = 479/700$ via $M_{22}$
\end{itemize}
Position the present paper's freeze-thaw transit as the dynamical face
of the substrate-algebraic structure derived in \cite{SandersGish2026M22}.
The Type-T/F/A taxonomy of \S\ref{sec:typeTFA} is the dynamical
expression of substrate paradox classes (forthcoming work).}

%%% ============================================================
\section{Summary}\label{sec:summary}

\TODO{One-paragraph wrap-up. Six key items as in existing draft, but
reframed:
\begin{enumerate}
\item Analytic vacuum $\Xi_0 = e^{-1}$ (geometric anchor of transit).
\item Stable mass $m_\Xi^2 = \Lambda^4 e / M_\Pl^2$;
$\Lambda \approx 1.7$~meV.
\item Dual-regime trajectory: T-thaw, F-frozen at $z \approx 2$,
A-asymptotic.
\item Structural connection $V = -\Lambda^4 H_\Gibbs$
(motivational).
\item No direct fifth force in minimal theory.
\item DESI consistency check at proximity level; full likelihood
in companion.
\end{enumerate}
Falsifiability via $F_1$--$F_5$ at Stage-IV precision.}

%%% ============================================================
\section*{Acknowledgments}

\TODO{Acknowledgments. Note verification scripts and DESI fit script
at zenodo deposit DOI \texttt{10.5281/zenodo.18852047}. All errors are
the authors'.}

\begin{thebibliography}{99}

%% ===== Foundational refs =====

\bibitem{BialynickiBirulaMycielski1976}
I.~Bialynicki-Birula and J.~Mycielski,
\emph{Nonlinear wave mechanics},
Ann. Phys. (N.Y.) \textbf{100}, 62--93 (1976).

\bibitem{RatraPeebles1988}
B.~Ratra and P.~J.~E.~Peebles,
\emph{Cosmological consequences of a rolling homogeneous scalar field},
Phys. Rev. D \textbf{37}, 3406 (1988).

\bibitem{Wetterich1988}
C.~Wetterich,
\emph{Cosmology and the fate of dilatation symmetry},
Nucl. Phys. B \textbf{302}, 668 (1988).

\bibitem{Frieman1995}
J.~A.~Frieman \emph{et al.},
\emph{Cosmology with ultralight pseudo-Nambu-Goldstone bosons},
Phys. Rev. Lett. \textbf{75}, 2077 (1995).

\bibitem{CPL2001}
M.~Chevallier and D.~Polarski,
\emph{Accelerating universes with scaling dark matter},
Int. J. Mod. Phys. D \textbf{10}, 213 (2001).

\bibitem{CaldwellLinder2005}
R.~R.~Caldwell and E.~V.~Linder,
\emph{The limits of quintessence},
Phys. Rev. Lett. \textbf{95}, 141301 (2005); arXiv:astro-ph/0505494.

\bibitem{BarrowParsons1995}
J.~D.~Barrow and P.~Parsons,
\emph{Inflationary models with logarithmic potentials},
Phys. Rev. D \textbf{52}, 5576 (1995); arXiv:astro-ph/9506049.

\bibitem{ColemanWeinberg1973}
S.~Coleman and E.~Weinberg,
Phys. Rev. D \textbf{7}, 1888 (1973).

%% ===== DESI =====

\bibitem{DESI2024BAO}
DESI Collaboration, arXiv:2404.03000 (2024).

\bibitem{DESI2024VI}
DESI Collaboration, arXiv:2404.03002 (2024).

\bibitem{DESI2025DR2}
DESI Collaboration, arXiv:2503.14738; Phys. Rev. D \textbf{112},
083515 (2025).

%% ===== Planck =====

\bibitem{Planck2020}
Planck Collaboration, A\&A \textbf{641}, A6 (2020).

%% ===== Stage-IV forecasts =====

\bibitem{ShajibFrieman2025}
A.~J.~Shajib and J.~A.~Frieman, arXiv:2502.06929 (2025).

\bibitem{Turyshev2026Review}
S.~G.~Turyshev, arXiv:2602.05368 (2026).

%% ===== Adjacent quintessence/log-DE literature =====

\bibitem{Thompson2019}
R.~I.~Thompson, MNRAS \textbf{482}, 5448 (2019); arXiv:1811.03164.

\bibitem{KouwnOh2012}
S.~Kouwn and P.~Oh, JKPS \textbf{65}, 814 (2014); arXiv:1201.4544.

\bibitem{Chavanis2015}
P.-H.~Chavanis, Phys. Rev. D \textbf{92}, 103004 (2015); arXiv:1505.00034.

\bibitem{Oikonomou2019}
V.~K.~Oikonomou, Ann. Phys. \textbf{409}, 167912 (2019); arXiv:1907.02600.

\bibitem{Wang2026QReconstruction}
S.~Wang \emph{et al.}, arXiv:2603.21125 (2026).

\bibitem{Adil2026Analytic}
S.~A.~Adil \emph{et al.}, arXiv:2603.14693 (2026).

\bibitem{AdamEtAl2025}
H.~Adam \emph{et al.}, arXiv:2509.13302 (2025).

%% ===== GW + local tests =====

\bibitem{GW170817}
LIGO/Virgo, ApJL \textbf{848}, L13 (2017).

\bibitem{EotWash}
D.~J.~Kapner \emph{et al.}, PRL \textbf{98}, 021101 (2007).

\bibitem{MICROSCOPE}
P.~Touboul \emph{et al.}, PRL \textbf{129}, 121102 (2022).

%% ===== Information-theoretic / entropic =====

\bibitem{Verlinde2011}
E.~Verlinde, JHEP \textbf{04}, 029 (2011).

\bibitem{CHLZ2012}
A.~Sheykhi \emph{et al.}, Astrophys. Space Sci. \textbf{342}, 93 (2012).

\bibitem{Ensslin2013}
T.~A.~En{\ss}lin, Ann. Phys. \textbf{525}, 895 (2013).

\bibitem{Caticha2012}
A.~Caticha, EBEB 2012.

%% ===== Companion papers =====

\bibitem{SandersGishSigma2026}
B.~R.~Sanders and M.~Gish,
\emph{Non-Associativity Decay in Binary Composition Tables over $\mathbb{Z}/N\mathbb{Z}$},
companion paper, in preparation (2026).

\bibitem{SandersGishFourCore2026}
B.~R.~Sanders and M.~Gish,
\emph{Joint Closure, Per-Coordinate Fuse Data, and a Closed-Form
Algebraic Attractor of Two Commutative Binary Operations on
$\mathbb{Z}/10\mathbb{Z}$},
companion paper, in preparation (2026).

\bibitem{SandersGish2026M22}
B.~R.~Sanders and M.~Gish,
\emph{Substrate-Algebraic Dark Energy: $\Omega_\DE = T^* - W/2$ from
Mathieu Group $M_{22}$ Mediation},
companion paper, in preparation (2026).

\end{thebibliography}

\end{document}
